\chapter{Лінійна алгебра та бібліотека NumPy.}

Лінійна алгебра — це розділ математики, що вивчає вектори, матриці та операції над ними. Вона є математичною основою більшості алгоритмів машинного навчання, аналізу даних, комп'ютерного зору та нейронних мереж. Методи лінійної алгебри використовуються для представлення даних у вигляді матриць, зменшення розмірності даних, побудови моделей машинного навчання, кластеризації, регресії, обробки зображень та аналізу текстів.

Вектори та матриці відіграють центральну роль в сучасному інтелектуальному аналізі даних. Матриці, це очевидний спосіб зберігання табличних даних. Сьогодні основа лінійної алгебри, є мовою всіх алгоритмів аналізу даних (машинне навчання, штучний інтелект та ін.).

\section{Основи лінійної алгебри.}

\textit{Скаляр} - окреме числове значення, яке характеризується лише величиною і не має напрямку.

\textit{Вектор} - 1D масив даних.

$\mathbf{v} = \begin{bmatrix} x_1 \\ x_2 \\ ... \\ x_n \end{bmatrix}$

\textit{Матриця} - 2D масив даних.

$$
\mathbf{A} = \begin{bmatrix}
x_{11} & x_{12} & ... & x_{1m} \\
x_{21} & x_{22} & ... & x_{2m} \\
... & ... & ... & ... \\
x_{n1} & x_{n2} & ... & x_{nm}
\end{bmatrix}
$$

Існують також узагальнення матриць вищого порядку (так звані *тензори*), які представляють тривимірні або вищі масиви значень. Вони досить широко використовуються в сучасній науці про дані, хоча зазвичай (але, звичайно, не завжди) тензори використовуються лише в значенні «багатовимірного масиву», а не в їхньому справжньому лінійно-алгебраїчному значенні. Тензори як лінійні оператори, що діють, наприклад, на матриці або інші тензори вищого порядку, дещо рідше використовуються в базовій науці про дані.


\subsection{Порядок рядків і стовпців.}

Матриці можуть бути розміщені в пам'яті за рядками або за стовпцями. Для того, щоб писати ефективний код для роботи із матричними даними, варто цей розуміти порядок запису.

$$
\mathbf{A} = \begin{bmatrix}
100 & 80 \\
60 & 80 \\
100 & 100
\end{bmatrix}
$$

За рядками: 100, 80, 60, 80, 100, 100. За стовпцями: 100, 60, 100, 80, 80, 100. 

\ding{43} По замовчуванню в NumPy порядок за рядками.



\subsection{Основні операції.}


\textbf{Додавання}: Дві матриці A та B, $[n\times m]$, $A+B=C$ визначається: $C_{ij}=A_{ij}+B_{ij}$

\textbf{Транспонування}: Для матриці $A$, $[n\times m]$, транспонована матриця $C=A^{T}$ визначається: $C_{ij}=A_{ji}$

\textbf{Множення}: Для матриць $A$ ($[n\times m]$) та $B$ ($[m\times p]$) добуток $C=AB$ ($[n\times p]$) визначається: $C_{ij}=\sum_{k=1}^{m}A_{ik}B_{kj}$

Множення матриць є асоціативним ($(AB)C=A(BC)$), дистрибутивним $A(B+C)=AB+AC$, але не комутативним ($AB\neq BA$).

\textbf{Одинична матриця}: Одинична матриця I це квадратна матриця з одиницями на головній діагоналі та нулями в інших позиціях:

$$
\mathbf{I} = \begin{bmatrix}
1 & 0 & ... & 0 \\
0 & 1 & ... & 0 \\
... & ... & ... & ... \\
0 & 0 & ... & 1
\end{bmatrix}
$$

Для будь-якої матриці $A$: $AI=IA=A$

\textbf{Обернена матриця}: Для квадратної матриці $A$, обернена матриця $A^{-1}$ це така матриця, що

$A^{-1}A=AA^{-1}=I$

Не для кожної квадратної матриці існує обернена.

Деякі рівняння: $(AB)^{T}=B^{T}A^{T}$, $(AB)^{-1}=B^{-1}A^{-1}$

Скалярний добуток векторів: $x*y$ = $x^{T}y = \sum_{i=1}^{n}x_{i}y_{i}$;

Норма вектора: $||x||_{2} = \sqrt{\sum_{i=1}^{n}x_{i}^{2}}$



\section{Масиви NumPy.}

Створення масиву за допомогою команди \verb|numpy.array| повертає тип \verb|ndarray|. Ви також можете створювати масиви нулів, одиниць або випадкових чисел (у цьому випадку \verb|np.randon.randn| створює матрицю зі стандартними випадковими нормальними елементами, а \verb|np.random.rand| створює рівномірні випадкові елементи).

Спочатку імпортуйте модуль NumPy за допомогою наступної команди:

\begin{lstlisting}
import numpy as np
\end{lstlisting}

Приклад створення вектора та масиву в NumPy:

\begin{lstlisting}
b = np.array([-13,9])           
A = np.array([[4,-5], [-2,3]])   
print(b, "\n")
print(A, "\n")
\end{lstlisting}

Результат:

\begin{console}
[-13   9] 

[[ 4 -5]
 [-2  3]] 
\end{console}

Приклад 1D масиву одиниць:

\begin{lstlisting}
print(np.ones(40), "\n")
\end{lstlisting}

Результат:

\begin{console}
[1. 1. 1. 1. 1. 1. 1. 1. 1. 1. 1. 1. 1. 1. 1. 1. 1. 1. 1. 1. 1. 1. 1. 1.
 1. 1. 1. 1. 1. 1. 1. 1. 1. 1. 1. 1. 1. 1. 1. 1.] 
\end{console}


1D масив нулів:

\begin{lstlisting}
print(np.zeros(4), "\n")
\end{lstlisting}

Результат:

\begin{console}
[0. 0. 0. 0.] 
\end{console}


1D масив сипадкових чисел нормального розподілу $N(0, 1)$:

\begin{lstlisting}
print(np.random.rand(40), "\n")
\end{lstlisting}

Результат:

\begin{console}
[0.18507018 0.46626434 0.9801947  0.0827781  0.57378447 0.56871301
 0.84269999 0.6657426  0.56345543 0.93262208 0.98343751 0.30900493
 0.38287176 0.01312199 0.23539386 0.15844975 0.90519996 0.26924736
 0.85669799 0.94183472 0.011728   0.63333182 0.80378177 0.64861789
 0.02308254 0.69805843 0.52183114 0.61711567 0.95150433 0.2085504
 0.76690385 0.1997735  0.45904879 0.81617043 0.62357921 0.8821446
 0.56465777 0.09903507 0.66269635 0.09785717]
\end{console}

Приклад 2D масиву одиниць ($3\times 4$, 3 рядки 4 стовпчики):

\begin{lstlisting}
print(np.ones((3,4)), "\n")
\end{lstlisting}

Результат:

\begin{console}
[[1. 1. 1. 1.]
 [1. 1. 1. 1.]
 [1. 1. 1. 1.]] 
\end{console}




2D масив нулів ($3\times 4$):

\begin{lstlisting}
print(np.zeros((3,4)), "\n")
\end{lstlisting}

Результат:

\begin{console}
[[0. 0. 0. 0.]
 [0. 0. 0. 0.]
 [0. 0. 0. 0.]] 
\end{console}



2D масив випадкових чисел нормального розподілу:

\begin{lstlisting}
print(np.random.randn(3,4)) 
\end{lstlisting}

Результат:

\begin{console}
[[-0.72206547  0.26200423  1.45715136 -1.25955877]
 [-0.01199544  0.27812632  2.34386974  0.1253467 ]
 [ 1.60420367  1.11023166  1.34434453  0.87136904]]
\end{console}

2D масив випадкових чисел нормального розподілу із середнім значенням $mu$ та дисперсією $sigma$:

\begin{lstlisting}
sigma, mu = 10, 50
print(sigma * np.random.randn(5, 5) + mu)
\end{lstlisting}

Результат:

\begin{console}
[[54.73972213 58.44438959 49.52225911 61.93038812 62.89356473]
 [58.20297378 54.47071892 53.323917   46.94830739 54.82418092]
 [61.0468657  30.48331679 57.33545183 84.93594736 55.92995469]
 [52.76434983 56.28182261 48.24665623 50.28029835 57.08573032]
 [44.30725611 48.51274777 40.88648182 36.13557207 55.81138676]]
\end{console}


Ви також можете створити матрицю ідентичності за допомогою команди \verb|np.eye()|, а діагональну матрицю — за допомогою команди \verb|np.diag()|.

Наступний приклад створює масив для одиничної матриці $3\times 3$, та діагональний масив випадкових чисел нормального розподілу:

\begin{lstlisting}
print(np.eye(3),"\n")
print(np.diag(np.random.randn(3)),"\n")
\end{lstlisting}

Результат:

\begin{console}
[[1. 0. 0.]
 [0. 1. 0.]
 [0. 0. 1.]] 

[[ 0.50651677  0.          0.        ]
 [ 0.         -0.35063699  0.        ]
 [ 0.          0.         -0.72949183]] 
\end{console}



\subsection{Індексація в масивах NumPy.}

В NumPy можна вибрати окремий елемент масиву, знаючи його положення в масиві, через індексацію \verb|A[i,j]|. Можна також вибрати цілий ряд \verb|A[i,:]|, або зріз масив (slice) \verb|A[start:end, :]|.

\begin{lstlisting}
A = np.array([[1,2,3], [4,5,6], [7,8,9], [10, 11, 12]])
print(A, "\n")
print(A[1,1],"\n")        
print(A[1,:],"\n")          
print(A[1:3, :], "\n")      
\end{lstlisting}

Результат:

\begin{console}
[[ 1  2  3]
 [ 4  5  6]
 [ 7  8  9]
 [10 11 12]] 

5 

[4 5 6] 

[[4 5 6]
 [7 8 9]] 
\end{console}

Можна також вибрати A[1,1] як одиничний 2D масив:

\begin{lstlisting}
print(A[1:2,1:2], "\n")
\end{lstlisting}

Результат:

\begin{console}
[[5]] 
\end{console}

Можна також вибрати окремі рядки, наприклад 1, 2, та 3:

\begin{lstlisting}
print(A[[1,2,3],:], "\n")
\end{lstlisting}

Результат:

\begin{console}
[[ 4  5  6]
 [ 7  8  9]
 [10 11 12]] 
\end{console}

Або, для прикладу, 2, 1, та 2 знову:

\begin{lstlisting}
print(A[[2,1,2],:], "\n")
\end{lstlisting}

Результат:

\begin{console}
[[7 8 9]
 [4 5 6]
 [7 8 9]] 
\end{console}

Можна також вибирати окремі рядки масиву, які нас цікавлять, вказуючи в окремому масиві через bool значення True або False, в залежності від того, чи ми хочемо вибрати відповідний рядок:

\begin{lstlisting}
print(A[[False, True, False, True],:], "\n")
\end{lstlisting}

Результат:

\begin{console}
[[ 4  5  6]
 [10 11 12]] 
\end{console}



\subsection{Основні операції над масивами в NumPy.}

Масиви можна додавати/віднімати, множити/ділити та транспонувати, але ці операції не є ідентичними до відповідних операцій у лінійній алгебрі.

Множення та ділення масивів виконуються поелементно, вони не є множенням матриць або чимось, пов'язаним з оберненням матриць.

Для прикладу, визначимо два масиви:

\begin{lstlisting}
A = np.array([[1,2,3], [4,5,6], [7,8,9], [10, 11, 12]])
B = np.array([[1, 1, 1], [1,2,1], [3, 1, 3], [1, 4, 1]])
\end{lstlisting}

Додати A та B поелементно:

\begin{lstlisting}
print(A+B, "\n")
\end{lstlisting}

Результат:

\begin{console}
[[ 2  3  4]
 [ 5  7  7]
 [10  9 12]
 [11 15 13]] 
\end{console}



Відняти B від A поелементно:

\begin{lstlisting}
print(A-B, "\n")
\end{lstlisting}

Результат:

\begin{console}
[[ 0  1  2]
 [ 3  3  5]
 [ 4  7  6]
 [ 9  7 11]] 
\end{console}



Поелементне множення (не множення матриць):

\begin{lstlisting}
print(A*B, "\n")
\end{lstlisting}

Результат:

\begin{console}
[[ 1  2  3]
 [ 4 10  6]
 [21  8 27]
 [10 44 12]] 
\end{console}



Поелементне ділення:

\begin{lstlisting}
print(A/B, "\n")
\end{lstlisting}

Результат:

\begin{console}
[[ 1.          2.          3.        ]
 [ 4.          2.5         6.        ]
 [ 2.33333333  8.          3.        ]
 [10.          2.75       12.        ]] 
\end{console}



Ви можете транспонувати масиви, але зверніть увагу, що це має значення тільки для 2D (або вищих) масивів. Транспонування 1D масиву нічого не дає, оскільки NumPy не розрізняє стовпцеві вектори та рядкові вектори для 1D масивів.

\begin{lstlisting}
print(A, "\n")
print(A.T, "\n")
\end{lstlisting}

Результат:

\begin{console}
[[ 1  2  3]
 [ 4  5  6]
 [ 7  8  9]
 [10 11 12]] 

[[ 1  4  7 10]
 [ 2  5  8 11]
 [ 3  6  9 12]] 
\end{console}




\subsection{Трансляція NumPy.}

Справжня цікавість починається, коли ви додаєте/віднімаєте/множите/ділите масиви різних розмірів. Замість того, щоб видавати помилку, NumPy спробує зрозуміти вашу операцію, використовуючи правила трансляції (broadcasting) NumPy. Це складна тема, яка часто бентежить новачків у NumPy, але з невеликою практикою правила стають досить інтуїтивними.

\begin{lstlisting}
A = np.ones((4,3))          # A - 4x3
x = np.array([[1,2,3]])      # x - 1x3

print(A, '\n')
print(x, '\n')
\end{lstlisting}

Результат:

\begin{console}
[[1. 1. 1.]
 [1. 1. 1.]
 [1. 1. 1.]
 [1. 1. 1.]] 

[[1 2 3]] 
\end{console}

Ми визначили масив $4\times 3$ та вектор розмірності $3$. 


\begin{lstlisting}
print(A*x)
\end{lstlisting}

Результат:

\begin{console}
[[1. 2. 3.]
 [1. 2. 3.]
 [1. 2. 3.]
 [1. 2. 3.]]
\end{console}

Коли ми викликали перемноження масиву та вектора, NumPy автоматично повторив $x$ 4 рази, та перемножив на масив.


\subsection{Операції лінійної алгебри в NumPy.}

Починаючи з Python 3, між масивами NumPy тепер існує оператор множення матриць \verb|@| (раніше для цього доводилося використовувати більш громіздку функцію \verb|np.dot()|).

\begin{lstlisting}
A = np.random.randn(5,4)
C = np.random.randn(4,3)
x = np.random.randn(4)
y = np.random.randn(5)
z = np.random.randn(4)
\end{lstlisting}

Приклад множення матриці на матрицю (вертає 2D масив):

\begin{lstlisting}
print(A @ C, "\n") 
\end{lstlisting}

Результат:

\begin{console}
[[-0.77983164  0.86915701  0.74016041]
 [ 0.9441811  -1.13337787 -2.45126146]
 [ 1.06065113 -3.17950771 -3.50084458]
 [ 1.26732128 -1.3041138  -2.5499486 ]
 [ 1.60891994 -4.68541868 -4.6433185 ]] 
\end{console}


Приклад множення матриці на вектор (вертає 1D масив):

\begin{lstlisting}
print(A @ x, "\n")
\end{lstlisting}

Результат:

\begin{console}
[ 1.53587182 -2.9122003  -2.53401694 -2.88846645 -3.50775768]
\end{console}


Внутрішній добуток (скалярний) - зверніть увагу на тип даних, що вертаються - це скаляр, а не масив!

\begin{lstlisting}
print(x @ z) 
\end{lstlisting}

Результат:

\begin{console}
2.236322572387176
\end{console}

\ding{43} За замовчуванням оператор \verb|@| застосовується зліва направо, що може призвести до дуже неефективних порядків множення матриць.

\begin{lstlisting}
A = np.random.randn(1000,1000)
B = np.random.randn(1000,2000)
x = np.random.randn(2000)

%timeit A @ B @ x
\end{lstlisting}

Результат:

\begin{console}
59.1 ms +/- 7.45 ms per loop (mean +/- std. dev. of 7 runs, 10 loops each)
\end{console}


Це виконує матричні добутки $(AB)x$, які спочатку обчислюють неефективне множення матриць. Якщо ми хочемо обчислити добуток у набагато ефективнішому (в даному прикладі) порядку $A(Bx)$, ми вказуємо порядок виконання операції множення матриць через дужки:


\begin{lstlisting}
%timeit A @ (B @ x)
\end{lstlisting}

Результат:

\begin{console}
2.13 ms +/- 226' ns per loop (mean +/- std. dev. of 7 runs, 100 loops each)
\end{console}


Нарешті, Numpy включає процедуру \verb|np.linalg.inv()| для обчислення оберненої матриці $A^{-1}$ та \verb|np.linalg.solve()| для обчислення матричного рівняння $A^{-1}b$.


\begin{lstlisting}
b = np.array([-13,9])
A = np.array([[4,-5], [-2,3]])

print(np.linalg.inv(A), "\n") 
print(np.linalg.solve(A,b))
\end{lstlisting}

Результат:

\begin{console}
[[1.5 2.5]
 [1.  2. ]] 

[3. 5.]
\end{console}





\subsection{Складність операцій.}

Припустимо, що $A$, $B$ [$n\times n$], $x,y$ - вектори [n].

Добуток матриць $AB$: $O(n^{3})$

Добуток матриці та вектора $Ax$: $O(n^{2})$

Внутрішній добуток векторів $x^{T}y$: $O(n)$

Обернена матриця/розв'язання: $A^{-1}$: $O(n^{3})$





\subsection{Розріджені матриці.}

Багато матриць є розрідженими - такими, що містять переважно нульові елементи, з лише декількома ненульовими 
елементами.

Приклади: матриці, утворені реальними графами, матриці підрахунку слів у документах 
(більше про це пізніше)

Зберігання всіх цих нулів у стандартному форматі матриці може бути великою тратою 
обчислювальних ресурсів і пам'яті.

Бібліотеки розріджених матриць надають ефективні засоби для обробки цих розріджених 
матриць, зберігаючи і оперуючи тільки ненульовими елементами

Існує кілька різних способів зберігання розріджених матриць, кожен з яких оптимізований для 
різних операцій

\textbf{Формат координат (COO):} зберігання кожного елемента як кортежу
($inde_{row}$, $index_{column}$, $value$)

Наприклад, замість матриці:

$$
\mathbf{A} = \begin{bmatrix}
0 & 0 & 3 & 0 \\
2 & 0 & 0 & 1 \\
0 & 1 & 0 & 0 \\
4 & 0 & 1 & 0
\end{bmatrix}
$$

Матимемо:

$$values = [2, 4, 1, 3, 1, 1]$$

$$indices_{row} = [1, 3, 2, 0, 3, 1]$$

$$indecies_{column} = [0, 0, 1, 2, 2, 3]$$

Приклад використання біблотеки \verb|scipy| для роботи з розрідженими матрицями:

\begin{lstlisting}
import scipy.sparse as sp

values = [2, 4, 1, 3, 1, 1]
row_indices = [1, 3, 2, 0, 3, 1]
column_indices = [0, 0, 1, 2, 2, 3]
A = sp.coo_matrix((values, (row_indices, column_indices)), shape=(4,4))
print(A.todense())
\end{lstlisting}

Результат:

\begin{console}
[[0 0 3 0]
 [2 0 0 1]
 [0 1 0 0]
 [4 0 1 0]]
\end{console}


\ding{46} \textbf{Питання для самоконтролю.}

\begin{enumerate}
\item Що таке лінійна алгебра і чому вона є важливою для інтелектуального аналізу даних?
\item Які задачі аналізу даних розв'язуються за допомогою методів лінійної алгебри?
\item Що таке скаляр?
\item Що таке вектор?
\item Що таке матриця?
\item Чим відрізняється вектор від матриці?
\item Які основні операції можна виконувати над векторами?
\item Які основні операції можна виконувати над матрицями?
\item Що таке транспонування матриці?
\item Що таке одинична матриця?
\item Що таке нульова матриця?
\item Що таке квадратна матриця?
\item Які умови повинні виконуватися для множення двох матриць?
\item Що таке скалярний добуток векторів?
\item Для чого використовується матричне множення в машинному навчанні?
\item Що таке визначник (детермінант) матриці?
\item Для чого використовується обернена матриця?
\item Що таке власні значення та власні вектори матриці?
\item Чому власні значення важливі для методів аналізу даних?
\item Що таке бібліотека NumPy?
\item Які переваги NumPy порівняно зі звичайними списками Python?
\item Як встановити бібліотеку NumPy?
\item Як імпортувати NumPy у Python-програму?
\item Що таке масив (ndarray) у NumPy?
\item Як створити одновимірний масив NumPy?
\item Як створити двовимірний масив NumPy?
\item Як отримати розмірність масиву?
\item Як визначити форму (shape) масиву?
\item Як створити масив, заповнений нулями?
\item Як створити масив, заповнений одиницями?
\item Для чого використовуються функції arange() та linspace()?
\item Як змінити форму масиву за допомогою reshape()?
\item Як звернутися до окремого елемента масиву?
\item Як виконати зрізи (slicing) масивів?
\item Як працює індексація багатовимірних масивів?
\item Як виконуються арифметичні операції над масивами NumPy?
\item Що таке векторизовані обчислення?
\item Які переваги векторизації над використанням циклів?
\item Як обчислити суму, середнє значення та стандартне відхилення елементів масиву?
\item Як знайти максимальне та мінімальне значення масиву?
\item Як виконати множення матриць у NumPy?
\item Для чого використовується функція dot()?
\item Як знайти транспоновану матрицю в NumPy?
\item Як обчислити обернену матрицю та визначник за допомогою NumPy?
\item Яку роль відіграють NumPy та лінійна алгебра в сучасному машинному навчанні та інтелектуальному аналізі даних?
\end{enumerate}
